\BourbakiExercises{习题}

\begin{enumerate}
\def\labelenumi{\arabic{enumi})}
\tightlist
\item
  设 B 是交换域 K 上的一个有限次数代数，b 是 B 的一个元素。
\end{enumerate}

\begin{enumerate}
\def\labelenumi{\alph{enumi})}
\item
  若 \(b\) 是幂零的，则对每个在 K 上有限维的 B-模 M，我们有 \(\mathrm{Tr}_{\mathrm{M / K}}(b) = 0\)。由此推出，若 \(b\) 属于 B 的根，则我们有 \(\mathrm{Tr}_{\mathrm{M / K}}(bc) = 0\)（对每个 \(c \in \mathrm{B}\)）。
\item
  证明：若 \(\operatorname{Tr}_{\mathrm{S}/\mathrm{K}}(bc)=0\) 对每个单 B-模 S 与每个 \(c\in B\) 成立，则 b 属于 B 的根（利用稠密性定理）。
\end{enumerate}

\begin{enumerate}
\def\labelenumi{\arabic{enumi})}
\setcounter{enumi}{1}
\tightlist
\item
  设 B 是交换域 K 上一个次数为 n 的有限代数。对 B 的任一元素 b，用 \(\mathrm{Pm}_{\mathrm{B}/\mathrm{K}}(b;\mathrm{X})\) 表示 b 在 K 上的最小多项式（V, §3, No.~1, 第 16 页）。我们有 \(\mathrm{Pm}_{\mathrm{B}/\mathrm{K}}(b;\mathrm{X})=\mathrm{Pm}_{\mathrm{B}^{\circ}/\mathrm{K}}(b;\mathrm{X})\)。
\end{enumerate}

\begin{enumerate}
\def\labelenumi{\alph{enumi})}
\item
  证明 \(\mathrm{Pm}_{\mathrm{B / K}}(b;\mathrm{X})\) 整除 \(\mathrm{Pc}_{\mathrm{B / K}}(b;\mathrm{X})\)，而后者整除 \((\mathrm{Pm}_{\mathrm{B / K}}(b;\mathrm{X}))^n\)。
\item
  对 K 的每个扩张 L，我们有 \(\mathrm{Pm}_{\mathrm{B}_{(\mathrm{L})}/\mathrm{L}}(1 \otimes b; \mathrm{X}) = \mathrm{Pm}_{\mathrm{B}/\mathrm{K}}(b; \mathrm{X})\)。
\item
  设 \((e_i)_{i\in \mathrm{I}}\) 是 B 在 K 上的一组基，\(\mathbf{Y} = (\mathrm{Y}_i)_{i\in \mathrm{I}}\) 是一族变量。代数 B 关于基 \((e_i)\) 的主多项式（记作 \(\mathrm{P}((e_i);\mathrm{X};\mathbf{Y})\)，或简记为 \(\mathrm{P}(\mathrm{X};\mathbf{Y})\)）是 \(\mathrm{K}(\mathbf{Y})\) 上元素 \(\sum_{i\in \mathrm{I}}\mathrm{Y}_i e_i\)（属于 \(\mathrm{B}_{(\mathrm{K}(\mathbf{Y}))}\)）的最小多项式。*证明它属于 \(\mathrm{K}[\mathrm{X};\mathbf{Y}]\)（参看 AC, VII, §3, n° 5, 第 230 页，引理 1 与定理 2）。
\item
  设 \((e_{i}^{\prime})\) 是 B 在 K 上的另一组基，\((\lambda_{ij})\) 是以 K 为元素的矩阵，使得 \(e_{i} = \sum_{j} \lambda_{ij} e_{j}^{\prime}\)（对 \(i \in I\)）；若令 \(Y_{i}^{\prime} = \sum_{j} \lambda_{ji} Y_{j}\)，则 \(\mathrm{P}((e_{i}); \mathrm{X}; \mathbf{Y}) = \mathrm{P}((e_{i}^{\prime}); \mathrm{X}; \mathbf{Y}^{\prime})\)。设 \(b = \sum_{i} \beta_{i} e_{i}\) 是 B 的一个元素；多项式 \(\mathrm{P}((e_{i}); \mathrm{X}; \beta_{1}, \ldots, \beta_{n})\) 与基 \((e_{i})\) 的选取无关；我们把它称为 b 在 K 上的主多项式。它整除 \(\mathrm{Pc}_{\mathrm{B/K}}(b; \mathrm{X})\)，且是 \(\mathrm{Pm}_{\mathrm{B/K}}(b; \mathrm{X}).*\) 的倍式。
\item
  设 \(B'\) 是另一个有限次数的 K-代数；对适当的基，计算代数 \(B \times B'\) 的主多项式。
\end{enumerate}

\begin{enumerate}
\def\labelenumi{\arabic{enumi})}
\setcounter{enumi}{2}
\tightlist
\item
  \begin{enumerate}
  \def\labelenumii{\alph{enumii})}
  \tightlist
  \item
    设 A 是 K 上的一个有限次数中心单代数。证明 A 的元素 \(a\) 的主多项式是 \(\mathrm{Pcrd}_{\mathrm{A} / \mathrm{K}}(a; \mathrm{X})\)。A 的主多项式（关于 A 在 K 上的一组基）就是 VIII, 第 345 页，引理 4 中所考虑的多项式 P（化归到 \(\mathrm{B} = \mathbf{M}_n(\mathrm{K})\) 的情形）。
  \end{enumerate}
\end{enumerate}

\begin{enumerate}
\def\labelenumi{\alph{enumi})}
\setcounter{enumi}{1}
\tightlist
\item
  设 B 是 K 上的一个绝对半单代数，\(\mathcal{S}\) 是单 B-模的类所成之集。证明 B 的元素 b 的主多项式等于多项式 \(\mathrm{Pc}_{\mathrm{S}/\mathrm{K}}(b;\mathrm{X})\)（S 取遍 \(\mathcal{S}\)）的乘积（化归到 K 代数闭的情形）。
\end{enumerate}

¶ 4) 设 D 是一个体，\(\sigma\) 是 D 的一个自同构，A 是环 D{[}X{]} \(_{\sigma}\)（VIII, 第 11 页），\(n\) 是一个 \(\geqslant 2\) 的整数，\(a\) 是 D 的一个非零元素。把 A 的元素 X \(^{n}\) - \(a\) 记作 \(f\)。\(a)\) 理想 A \(f\) 是双边的，当且仅当 \(\sigma(a) = a\) 且 \(\sigma^{n}(x) = axa^{-1}\)（对每个 \(x \in \mathrm{D}\)）。

\begin{enumerate}
\def\labelenumi{\alph{enumi})}
\setcounter{enumi}{1}
\tightlist
\item
  假设包含在 D 的中心中的 n 次单位根所成的群 \(\mu_{n}\) 是 n 阶循环群且被 \(\sigma\) 固定。设 \(g \in A\) 是一个右除 f 的不可约首一多项式（VIII, 第 22 页，习题 27）。证明 g 的次数 d 整除 n，且 f 是 n/d 个次数为 d 的多项式的乘积。（设 h 是使 \(Ah = \bigcap_{\zeta \in \mu_{n}} Ag_{\zeta}\) 成立的首一多项式，其中 \(g_{\zeta}(X) = g(\zeta X)\)。注意我们必然有 \(h(\zeta X) = h(X)\)（对每个 \(\zeta \in \mu_{n}\)），并由此推出 h = f。然后，对 \(1 \leqslant i \leqslant n/d\)，构造一个元素 \(\zeta_{i}\)（属于 \(\mu_{n}\)）与一个次数为 id 的多项式 \(h_{i}\)，使得 \(Ag_{\zeta_{1}} \cap \cdots \cap Ag_{\zeta_{i}} = Ah_{i}\)，从而应用同引文 b)。
\end{enumerate}

此外，若理想 Af 是双边的，则环 A/Af 是半单的，且其单分量都有相同的长度（注意每个单 \((\mathrm{A}/\mathrm{Af})\)-模同构于模 \(A/Ag_{\zeta}\) 之一）。

\begin{enumerate}
\def\labelenumi{\alph{enumi})}
\setcounter{enumi}{2}
\tightlist
\item
  从现在起假设 b) 的假设成立、理想 A f 是双边的，且 n 是一个素数。由 b) 推出我们处于下述情形之一：
\end{enumerate}

\begin{enumerate}
\def\labelenumi{(\roman{enumi})}
\item
  存在一个 \(b \in D\)，使得 \(a = b^{n}\)，且 \(\sigma\) 是与 b 相伴的内自同构。环 C = A/Af 同构于积环 \(D^{n}\)。
\item
  (i) 的条件不满足，但存在一个 \(c \in \mathrm{D}\)，使得 \(a = \sigma^{n-1}(c)\sigma^{n-2}(c)\ldots\sigma(c)c\)。环 C 是长度为 \(n\) 的单环。
\item
  不存在 D 的元素 c 使得 \(a = \sigma^{n-1}(c) \sigma^{n-2}(c) \cdots \sigma(c) c\)。环 C 是一个体。
\end{enumerate}

在 D 交换的情形，重新得出 VIII, 第 330 页，习题 11 的结果。此外，若 n = 2，则 C 是 D 的在 \(\sigma\) 下不变的子域 K 上的一个四元数代数。

\begin{enumerate}
\def\labelenumi{\alph{enumi})}
\setcounter{enumi}{3}
\item
  证明存在 C 的一个自同构 \(\tau\)，其阶为 \(n\)，且其不变元子域等于 D。在上面的情形 (ii) 与 (iii) 中，D 是 C 的一个伽罗瓦子域（VIII, 第 274 页，习题 16）。
\item
  假设 \(\sigma\) 是与 D 的元素 \(d\) 相伴的内自同构。证明我们有 \(a = d^n z\)，其中 \(z\) 属于 C 的中心 Z，且环 C 同构于张量积 \(\mathrm{D} \otimes_{\mathbb{Z}} \mathrm{Z}[\mathrm{X}] / (\mathrm{X}^n - z)\)。
\item
  假设 \(\sigma\) 不是内自同构。证明 C 的中心是由 Z 中在 \(\sigma\) 下不变的元素所组成的子域。
\end{enumerate}

¶ 5) 设 \(K_{0}\) 是有理分式域 \(\mathbf{Q}(\mathrm{X})\)。设 K 是二次扩张 \(\mathrm{K}_{0}(\rho)\)，其中 \(\rho\) 是一个本原三次单位根，D 是 K 的循环扩张 \(\mathrm{K}(\sqrt[3]{\mathrm{X}})\)，\(\sigma\) 是 D 在 K 上的伽罗瓦群的一个生成元。考虑如习题 4 中那样定义的代数 C，取 n = 3 且 a = 2。

\begin{enumerate}
\def\labelenumi{\alph{enumi})}
\item
  证明 C 是一个以 K 为中心的体，且在 K 上的次数为 9。（要看出 D 没有范数为 2 的元素，化为证明不可能有形如 \(2F^{3} = P^{3} + XQ^{3} + X^{2}R^{3} + 3XPQR\) 的关系，其中 F, P, Q, R 是 \(\mathbf{Q}(\rho)[\mathrm{X}]\) 中两两互素的多项式。为此，考察这些多项式的最低次项。）
\item
  证明不存在 C 的自同构延拓 K 的异于恒同映射的唯一的 \(K_{0}\)-自同构（化归到这样的自同构保持 D 的情形）。由此推出 \(K_{0}\) 不是 C 的伽罗瓦子域（VIII, 第 274 页，习题 16），且不存在 C 的阶为 2 的 \(K_{0}\)-自同构。
\item
  设 \(\xi\) 是 D 的一个使 \(\xi^3 = \mathrm{X}\) 的元素，\(\eta\) 是 C 的一个使 \(\eta^3 = 2\) 且 \(\eta \xi = \rho \xi \eta\) 的元素；令 \(\zeta = \xi + \eta + \eta^2\)。证明 \(\mathrm{K}(\zeta)\) 是 D 的一个极大交换子域且在 K 上不是伽罗瓦的，而极大交换子域 \(\mathrm{K}(\xi)\) 与 \(\mathrm{K}(\eta)\) 在 K 上是伽罗瓦的（计算 \(\zeta^3\)）。
\end{enumerate}

\begin{enumerate}
\def\labelenumi{\arabic{enumi})}
\setcounter{enumi}{5}
\tightlist
\item
  设 D 是一个以 K 为中心、在 K 上的次数为 \(m^{2}\) 的体。设 x 是 D 的一个在 K 上次数为 n（\textgreater{} 1）的元素，\(f(\mathrm{X})\) 是它在 K 上的最小多项式。
\end{enumerate}

\begin{enumerate}
\def\labelenumi{\alph{enumi})}
\item
  证明 D{[}X{]}（VIII, 第 11 页）中一个次数 \textless{} n 的首一多项式不可能被所有多项式 \(X - txt^{-1}\)（\(t \in D^{*}\)）右整除（注意理想 \(\cap(\mathrm{D}[\mathrm{X}](\mathrm{X}-txt^{-1}))\) 的首一生成元的系数属于 K）。
\item
  设 \(u(\mathrm{X}), v(\mathrm{X}), w(\mathrm{X})\) 是 D{[}X{]} 中三个首一多项式，满足 \(u(\mathrm{X}) = v(\mathrm{X})w(\mathrm{X})\)，且 \(\mathrm{X} - x\) 是 \(u(\mathrm{X})\) 的右因子而不是 \(w(\mathrm{X})\) 的右因子。写 \(w(\mathrm{X}) = q(\mathrm{X})(\mathrm{X} - x) + t\)，其中 \(t \in \mathrm{D}^*\)。证明 \(\mathrm{X} - txt^{-1}\) 是 \(v(\mathrm{X})\) 的右因子。
\item
  由 a) 与 b) 推出，存在 D 的 n 个元素 \(t_{1},\ldots,t_{n}\)，使得我们有
\end{enumerate}

\[
f (\mathrm{X}) = \left(\mathrm{X} - t _ {1} x t _ {1} ^ {- 1}\right) \dots \left(\mathrm{X} - t _ {n} x t _ {n} ^ {- 1}\right)
\]

（注意 \(f(\mathrm{X})\) 被所有多项式 \(\mathrm{X} - txt^{-1}\) 右整除）。

\begin{enumerate}
\def\labelenumi{\alph{enumi})}
\setcounter{enumi}{3}
\tightlist
\item
  由此推出，\(\operatorname{Trd}_{\mathrm{D}/\mathrm{K}}(x)\)（相应地，\(\operatorname{Nrd}_{\mathrm{D}/\mathrm{K}}(x)\)）是 m 个形如 \(txt^{-1}\) 的元素之和（相应地，之积）（应用 VIII, 第 347 页，命题 8）。
\end{enumerate}

\begin{enumerate}
\def\labelenumi{\arabic{enumi})}
\setcounter{enumi}{6}
\tightlist
\item
  设 D 是一个以 K 为中心、在 K 上为有限秩的体，带有一个与其加群结构相容的全序，且两个正元素之积为正。证明 D 等于 K。（证明 K 的特征为零。若 \(D \neq K\)，注意在 D-K 中存在使 \(\operatorname{Trd}_{\mathrm{D}/\mathrm{K}}(x) = 0\) 的元素 x，并应用习题 6, d)）。
\end{enumerate}

